\documentclass[UTF8]{ctexart}
\usepackage{geometry,booktabs,longtable,array,tabularx}
\geometry{a4paper,margin=2.0cm}
\setlength{\emergencystretch}{3em}
\setlength{\tabcolsep}{3pt}
\renewcommand{\arraystretch}{1.12}
\title{JiuwenSwarm v4 可审计科研论文半流程报告}
\author{JiuwenSwarm/UCR}
\date{}
\begin{document}
\maketitle
\begin{abstract}
本文汇总三个主题的科研半流程演示。系统检索公开资料元数据，执行 JiuwenSwarm/UCR 本地证据实验，生成研究计划和报告，并保存 Claim--Evidence 账本。本文不把流程基准包装成领域科学结论；正式引用、实验真实性、研究结论和最终提交资格仍需人工确认。
\end{abstract}
\section{统一边界}
UCR 的标签是 \texttt{process\_evidence\_only}，它表示系统的流程证据和审计能力，不等于上下文工程、记忆引擎或自我进化的领域效果。所有来源默认处于待审核状态。当前 PDF 是由本机 XeLaTeX 编译生成的正式查看产物，源文件和实验记录同时保留在提交包中。
\section{矩阵概览}
\begin{longtable}{@{}>{\raggedright\arraybackslash}p{0.30\textwidth}>{\raggedright\arraybackslash}p{0.42\textwidth}>{\raggedleft\arraybackslash}p{0.14\textwidth}@{}}
\toprule
主题 & 状态 & Provider 调用\\\\
\midrule
\endfirsthead
\toprule
主题 & 状态 & Provider 调用\\\\
\midrule
\endhead
context-\allowbreak{}engineering & \texttt{\scriptsize completed\_\allowbreak{}pending\_\allowbreak{}human\_\allowbreak{}review} & 0 \\
memory-\allowbreak{}engine & \texttt{\scriptsize completed\_\allowbreak{}pending\_\allowbreak{}human\_\allowbreak{}review} & 0 \\
self-\allowbreak{}evolution & \texttt{\scriptsize completed\_\allowbreak{}pending\_\allowbreak{}human\_\allowbreak{}review} & 0 \\
\bottomrule
\end{longtable}
\section{Agent 上下文工程}
\textbf{摘要：}本文展示可审计的 JiuwenSwarm / UCR 流程基准，不把流程结果包装成领域科学结论。\\
\textbf{结论：}v4 形成可重复、可审计的研究半流程闭环；最终研究结论仍需人工确认。\\
\textbf{限制：}资料尚未完成人工引用审批；当前实验是流程证据基准，不足以证明三个主题的领域效果。
\section{Agent 记忆引擎}
\textbf{摘要：}本文展示可审计的 JiuwenSwarm / UCR 流程基准，不把流程结果包装成领域科学结论。\\
\textbf{结论：}v4 形成可重复、可审计的研究半流程闭环；最终研究结论仍需人工确认。\\
\textbf{限制：}资料尚未完成人工引用审批；当前实验是流程证据基准，不足以证明三个主题的领域效果。
\section{Agent 自我进化}
\textbf{摘要：}本文展示可审计的 JiuwenSwarm / UCR 流程基准，不把流程结果包装成领域科学结论。\\
\textbf{结论：}v4 形成可重复、可审计的研究半流程闭环；最终研究结论仍需人工确认。\\
\textbf{限制：}资料尚未完成人工引用审批；当前实验是流程证据基准，不足以证明三个主题的领域效果。
\section{编译与审核说明}
该统一报告和三个主题报告都需要人工阅读。编译命令为：
\begin{verbatim}
xelatex -interaction=nonstopmode -halt-on-error paper.tex
xelatex -interaction=nonstopmode -halt-on-error paper.tex
\end{verbatim}
最终状态保持为 \texttt{prepared\_not\_uploaded}。
\end{document}
